\chapter{Course Information}

\section*{Course materials}

We will use these course notes and \emph{Number Theory} by Andrews \cite{andrews1994}.

\section*{Important dates}
\begin{itemize}
\item First lecture: January 21st.
\item First midterm: February 25th.
\item Spring break: March 17-24th.
\item Second midterm: April 1st.
\item Final lecture: April 24th
\item Final exam: May 8th.
\end{itemize}

The first midterm covers content up to February 13th. The second midterm covers content up to the spring break.

\begin{enumerate}
\item Midterm $1$ covers \cref{chapter:induction,chapter:euclidean-division-and-representability,chapter:linear-diophantine-equations,chapter:fundamental-theorem-of-arithmetic,chapter:combinatorial-methods} from these notes, and Chapters 1, 2, and 3 from Andrews \emph{Number Theory}. Note that parts of Chapter 3 of Andrews' book do not have an equivalent in these notes, so don't forget to study Chapter 3 of Andrews.
\item Midterm $2$ covers \cref{chapter:polynomials,chapter:congruences,chapter:modular-arithmetic,chapter:systems-of-linear-congruences,chapter:polynomial-congruences,chapter:primitive-roots} from these notes, and Chapters 4, 5, 6.1, and 7 from Andrews \emph{Number Theory}.
\end{enumerate}

\newpage

\section*{Grading}
The course grade will be determined as follows:

\begin{itemize}
    \item Homework: 60\%
    \item Each exam (two midterms and final): 20\% each
\end{itemize}

This adds up to a total of 120\%, allowing students the opportunity to earn extra credit. The grading scale is as follows:

\begin{itemize}
    \item $>100\%$: A+
    \item $90$--$100\%$: A
    \item $80$--$90\%$: B
    \item $70$--$80\%$: C
    \item $60$--$70\%$: D
    \item $<60\%$: F
\end{itemize}

If the median score for the class falls below a B, grades may be adjusted (curved) so that the median corresponds to a B. Any adjustments will be made consistently for all students.

%%% Local Variables:
%%% mode: latex
%%% TeX-master: "lectures-on-number-theory"
%%% End:
